\subsection{非线性算子理论方法}\label{sec:nonlinear-operator-theory-methods}

本节给出从 Banach 空间 $X$ 到其对偶空间的非线性映射的一些结果:
\begin{equation}\label{eq:nonlinear-banach-to-dual-mapping}
X\ni v\mapsto T(v)\in X^*.
\end{equation}
我们通常用\emph{映射}一词指非线性映射, 而将\emph{算子}一词留给假定为线性的映射. 不过, 读者须注意, 这只是本书采用的一种人为约定, 并非全部文献都这样使用. 作者们既用 ``算子\textquotesingle\textquotesingle\ 指线性映射, 也用它指非线性映射; 在只研究线性算子的语境中, 常常不明确说明算子的线性假设. (Chapter~\ref{chap:operator-theory} 就采用这种做法.)

\subsubsection{有限维空间上的映射}\label{subsec:finite-dimensional-nonlinear-mappings}

本节研究映射 $\bm{f}:\mathbb{R}^n\to\mathbb{R}^n$. 目标是更好地理解 ``$n$ 个未知量的 $n$ 个方程\textquotesingle\textquotesingle\ 这一问题:
\begin{equation}\label{eq:finite-dimensional-nonlinear-system}
\bm{f}(\bm{u})=\bm{a}.
\end{equation}
在 Chapter~\ref{chap:introduction} 中, 已经讨论了这一问题的一个 ``局部\textquotesingle\textquotesingle\ 结果, 即逆函数定理; 但我们希望得到更强的结果. 作为所期待结果的一个例子, 考虑 $n=1$ 时的下面这一结果.

\begin{Theorem}
\label{thm:scalar-continuous-monotone-surjectivity}
假设 $f:\mathbb{R}\to\mathbb{R}$ 连续, 且满足
\begin{equation}\label{eq:scalar-continuous-surjectivity-growth-condition}
f(u)\to\pm\infty\quad\text{as }u\to\pm\infty.
\end{equation}
则对每个 $a\in\mathbb{R}$, 方程
\begin{equation}\label{eq:scalar-continuous-surjectivity-equation}
f(u)=a
\end{equation}
有解 $u\in\mathbb{R}$. 此外, 若 $f$ 单调递增, 则对每个 $a$, 解集构成一个闭区间; 若 $f$ 严格单调, 则解唯一.
\end{Theorem}

证明是介值定理以及单调函数和连续函数定义的一个初等应用.

为对从 $\mathbb{R}^n$ 到 $\mathbb{R}^n$ 的映射得到类似结果, 必须同时推广增长条件 \eqref{eq:scalar-continuous-surjectivity-growth-condition} 和单调递增 (保序) 的定义.

\begin{Definition}
\label{def:finite-dimensional-coercive-mapping}
若
\begin{equation}\label{eq:finite-dimensional-nonlinear-coercivity-definition}
\dfrac{\bm{f}(\bm{u})\cdot\bm{u}}{|\bm{u}|}\to\infty\quad\text{as }|\bm{u}|\to\infty.
\end{equation}
则称函数 $\bm{f}:\mathbb{R}^n\to\mathbb{R}^n$ 为\textbf{强制的}.
\end{Definition}

\begin{Definition}
\label{def:finite-dimensional-monotone-mapping}
若
\begin{equation}\label{eq:finite-dimensional-monotonicity-definition}
(\bm{f}(\bm{u})-\bm{f}(\bm{v}))\cdot(\bm{u}-\bm{v})\geq0\quad\text{for all }\bm{u},\bm{v}\in\mathbb{R}^n.
\end{equation}
则称函数 $\bm{f}:\mathbb{R}^n\to\mathbb{R}^n$ 为单调的. 若每当 $\bm{u}\neq\bm{v}$ 时, \eqref{eq:finite-dimensional-monotonicity-definition} 中的不等式都严格成立, 则称 $\bm{f}$ 为严格单调的.
\end{Definition}

首先研究单调性的推论. 从定义的下面这一直接推论开始.

\begin{Theorem}
\label{thm:finite-dimensional-strict-monotonicity-uniqueness}
若 $\bm{f}$ 严格单调, 则方程
\begin{equation}\label{eq:finite-dimensional-strictly-monotone-equation}
\bm{f}(\bm{u})=\bm{a}
\end{equation}
至多有一个解.
\end{Theorem}

为在 $\bm{f}$ 仅为单调时得到结果, 先证明下面的引理.

\begin{Lemma}
\label{lem:finite-dimensional-minty-equivalence}
给定 $\bm{a}\in\mathbb{R}^n$, 并设 $\bm{f}:\mathbb{R}^n\to\mathbb{R}^n$ 连续. 若 $\bm{u}\in\mathbb{R}^n$ 为变分不等式
\begin{equation}\label{eq:finite-dimensional-minty-variational-inequality}
(\bm{f}(\bm{v})-\bm{a})\cdot(\bm{v}-\bm{u})\geq0\quad\text{for every }\bm{v}\in\mathbb{R}^n,
\end{equation}
的解, 则 $\bm{u}$ 也是方程 \eqref{eq:finite-dimensional-strictly-monotone-equation} 的解.

此外, 若 $\bm{f}$ 单调, 则方程 \eqref{eq:finite-dimensional-strictly-monotone-equation} 的每个解也都是变分不等式 \eqref{eq:finite-dimensional-minty-variational-inequality} 的解.
\end{Lemma}

\begingroup
\renewcommand{\endProof}{\par\nobreak\hfill$\blacksquare$\par\vspace{0.1cm}}
\begin{Proof}
假设 $\bm{f}$ 连续, 且 \eqref{eq:finite-dimensional-minty-variational-inequality} 成立. 对任意 $\bm{w}\in\mathbb{R}^n$, 当 $t>0$ 时令 $\bm{v}=\bm{u}+t\bm{w}$. 由此得到 (将 \eqref{eq:finite-dimensional-minty-variational-inequality} 除以 $t$ 后)
\begin{equation}\label{eq:finite-dimensional-minty-directional-inequality}
(\bm{f}(\bm{u}+t\bm{w})-\bm{a})\cdot\bm{w}\geq0\quad\text{for all }\bm{w}\in\mathbb{R}^n.
\end{equation}
令 $t\to0$, 得到
\begin{equation}\label{eq:finite-dimensional-minty-directional-limit}
(\bm{f}(\bm{u})-\bm{a})\cdot\bm{w}\geq0\quad\text{for all }\bm{w}\in\mathbb{R}^n.
\end{equation}
取 $\bm{w}=\pm(\bm{f}(\bm{u})-\bm{a})$, 即得 $\bm{f}(\bm{u})=\bm{a}$.

现在假设 $\bm{f}$ 单调, 且 \eqref{eq:finite-dimensional-strictly-monotone-equation} 成立. 则由单调性的定义有
\begin{equation}\label{eq:finite-dimensional-minty-monotonicity-converse}
(\bm{f}(\bm{v})-\bm{a})\cdot(\bm{v}-\bm{u})=(\bm{f}(\bm{v})-\bm{f}(\bm{u}))\cdot(\bm{v}-\bm{u})\geq0\quad\text{for every }\bm{v}\in\mathbb{R}^n,
\end{equation}

\end{Proof}
\endgroup

作为这一结果的直接推论, 得到下面的结论.

\begin{Theorem}
\label{thm:finite-dimensional-monotone-solution-set-closed-convex}
设 $\bm{a}\in\mathbb{R}^n$, 并设 $\bm{f}$ 连续且单调. 则 $\bm{f}(\bm{u})=\bm{a}$ 的解集 $\mathcal{K}$ 为闭凸集.
\end{Theorem}

\begingroup
\renewcommand{\endProof}{\par\nobreak\hfill$\blacksquare$\par\vspace{0.1cm}}
\begin{Proof}
对任意固定的 $\bm{v}\in\mathbb{R}^n$, 令
\begin{equation}\label{eq:monotone-solution-set-halfspaces}
S_{\bm{v}}=\{\bm{u}\in\mathbb{R}^n\mid(\bm{f}(\bm{v})-\bm{a})\cdot(\bm{v}-\bm{u})\geq0\}.
\end{equation}
注意, 每个集合 $S_{\bm{v}}$ 都是闭半空间 (因而是凸的). 由 Lemma~\ref{lem:finite-dimensional-minty-equivalence}, 解集可写成
\begin{equation}\label{eq:monotone-solution-set-halfspace-intersection}
\mathcal{K}=\bigcap_{\bm{v}\in\mathbb{R}^n}S_{\bm{v}}.
\end{equation}
由于闭集的任意交为闭集, 凸集的任意交为凸集, 定理得证.
\end{Proof}
\endgroup

现在考察强制性的推论.

\begin{Theorem}
\label{thm:finite-dimensional-coercive-surjectivity}
设 $\bm{f}:\mathbb{R}^n\to\mathbb{R}^n$ 连续且强制. 则对每个 $\bm{a}\in\mathbb{R}^n$, 方程 \eqref{eq:finite-dimensional-strictly-monotone-equation} 都有解 $\bm{u}\in\mathbb{R}^n$.
\end{Theorem}

为证明这一结论, 需要下面的重要结果.

\begin{Theorem}{Brouwer 不动点定理}
\label{thm:brouwer-fixed-point}
设 $C$ 为 $\mathbb{R}^n$ 的非空紧凸子集, 并假设 $\bm{f}$ 为将 $C$ 映入 $C$ 的连续函数. 则 $\bm{f}$ 在 $C$ 中有不动点; 即存在 $\bm{u}\in C$, 使
\begin{equation}\label{eq:brouwer-fixed-point-equation}
\bm{f}(\bm{u})=\bm{u}.
\end{equation}

\end{Theorem}

\begingroup
\renewcommand{\endProof}{\par\nobreak\hfill$\blacksquare$\par\vspace{0.1cm}}
\begin{Proof}
一般情形请读者参阅文献 (见 \cite{DS}); 但当 $n=1$ 时, 证明容易从图形上看出. 此时定理的假设是, 连续函数 $f$ 将区间 $[a,b]$ 映入自身. 函数图像必位于 $(x,y)$ 平面内的方框 $[a,b]\times[a,b]$ 中. 只有三种可能: $f(a)=a$, $f(b)=b$, 或 $f$ 的图像从直线 $x=y$ 的左侧开始, 在其右侧结束. 由于函数连续, 图像必在某处穿过直线 $x=y$. 任一交点都给出一个不动点.
\end{Proof}
\endgroup

现在证明 Theorem~\ref{thm:finite-dimensional-coercive-surjectivity}.

\begingroup
\renewcommand{\endProof}{\par\nobreak\hfill$\blacksquare$\par\vspace{0.1cm}}
\begin{Proof}
首先注意到, 只需证明 $\bm{f}(\bm{u})=0$ 有解. 为看出这一点, 任给 $\bm{a}\in\mathbb{R}^n$, 定义
\begin{equation}\label{eq:coercive-surjectivity-shifted-mapping}
\bm{f}_{\bm{a}}(\bm{u}):=\bm{f}(\bm{u})-\bm{a}.
\end{equation}
注意, $\bm{f}_{\bm{a}}$ 为强制的当且仅当 $\bm{f}$ 为强制的, 因为\[\begin{aligned}
\dfrac{\bm{f}_{\bm{a}}(\bm{u})\cdot\bm{u}}{|\bm{u}|}&=\dfrac{\bm{f}(\bm{u})\cdot\bm{u}}{|\bm{u}|}-\dfrac{\bm{a}\cdot\bm{u}}{|\bm{u}|}\\
&\geq\dfrac{\bm{f}(\bm{u})\cdot\bm{u}}{|\bm{u}|}-\dfrac{|\bm{a}||\bm{u}|}{|\bm{u}|}\\
&=\dfrac{\bm{f}(\bm{u})\cdot\bm{u}}{|\bm{u}|}-|\bm{a}|.
\end{aligned}\]因此, 若能证明对每个强制函数, $\bm{f}(\bm{u})=0$ 都有解, 就立即得到对每个 $\bm{a}\in\mathbb{R}^n$, $\bm{f}(\bm{u})=\bm{a}$ 都有解.

现在将 $\bm{f}(\bm{u})=0$ 转化为不动点问题. 第一步, 定义 $\bm{g}:\mathbb{R}^n\to\mathbb{R}^n$ 为
\begin{equation}\label{eq:coercive-surjectivity-fixed-point-mapping}
\bm{g}(\bm{u}):=\bm{u}-\bm{f}(\bm{u}).
\end{equation}
注意到
\begin{equation}\label{eq:coercive-surjectivity-fixed-point-pairing}
\bm{g}(\bm{u})\cdot\bm{u}=|\bm{u}|^2-\bm{f}(\bm{u})\cdot\bm{u}.
\end{equation}
由于 $\bm{f}$ 为强制的, 存在 $R>0$, 使得
\begin{equation}\label{eq:coercive-surjectivity-large-radius-strict-bound}
\bm{g}(\bm{u})\cdot\bm{u}<|\bm{u}|^2\quad\text{whenever }|\bm{u}|\geq R.
\end{equation}
现在定义
\begin{equation}\label{eq:coercive-surjectivity-ball-retraction}
\bm{r}(\bm{v}):=\begin{cases}\bm{v},&|\bm{v}|\leq R,\\\dfrac{R\bm{v}}{|\bm{v}|},&|\bm{v}|>R.\end{cases}
\end{equation}
最后令
\begin{equation}\label{eq:coercive-surjectivity-retracted-fixed-point-mapping}
\bm{h}(\bm{u}):=\bm{r}(\bm{g}(\bm{u})).
\end{equation}
注意, $\bm{h}$ 为连续函数的复合, 因而连续. 此外, 若令
\begin{equation}\label{eq:coercive-surjectivity-open-ball}
B:=\{\bm{u}\in\mathbb{R}^n\mid|\bm{u}|<R\},
\end{equation}
则 $\bm{h}$ 将闭球 $\overline B$ 映入自身. 因此, 由 Brouwer 不动点定理, $\bm{h}$ 有不动点 $\bm{u}\in\overline B$.

若 $\bm{u}\in\partial B$, 则
\begin{equation}\label{eq:coercive-surjectivity-boundary-fixed-point-norms}
R=|\bm{u}|=|\bm{h}(\bm{u})|=|\bm{r}(\bm{g}(\bm{u}))|.
\end{equation}
因此, 由 $\bm{r}$ 的定义必有
\begin{equation}\label{eq:coercive-surjectivity-boundary-radius-ratio}
\dfrac{R}{|\bm{g}(\bm{u})|}\leq1.
\end{equation}
但现在利用 \eqref{eq:coercive-surjectivity-large-radius-strict-bound}, 得到
\begin{equation}\label{eq:coercive-surjectivity-boundary-contradiction}
|\bm{u}|^2=\bm{u}\cdot\bm{u}=\bm{h}(\bm{u})\cdot\bm{u}=\dfrac{R}{|\bm{g}(\bm{u})|}\bm{g}(\bm{u})\cdot\bm{u}<\dfrac{R}{|\bm{g}(\bm{u})|}|\bm{u}|^2\leq|\bm{u}|^2.
\end{equation}
这构成矛盾. 因此必有
\begin{equation}\label{eq:coercive-surjectivity-interior-fixed-point-norms}
R>|\bm{u}|=|\bm{h}(\bm{u})|=|\bm{r}(\bm{g}(\bm{u}))|.
\end{equation}
从而 $\bm{r}(\bm{g}(\bm{u}))=\bm{g}(\bm{u})$, 故
\begin{equation}\label{eq:coercive-surjectivity-interior-retraction-identity}
\bm{u}=\bm{h}(\bm{u})=\bm{r}(\bm{g}(\bm{u}))=\bm{g}(\bm{u})=\bm{u}-\bm{f}(\bm{u}).
\end{equation}
最后得到
\begin{equation}\label{eq:coercive-surjectivity-zero-equation-conclusion}
\bm{f}(\bm{u})=0,
\end{equation}
证明完成.
\end{Proof}
\endgroup

综合本节的结果, 得到

\begin{Corollary}
\label{cor:finite-dimensional-coercive-strict-monotone-bijection}
若 $\bm{f}:\mathbb{R}^n\to\mathbb{R}^n$ 连续, 严格单调且强制, 则对每个 $\bm{a}\in\mathbb{R}^n$, 存在唯一的 $\bm{u}\in\mathbb{R}^n$, 使\[\bm{f}(\bm{u})=\bm{a}.\]
\end{Corollary}

\subsubsection{Banach 空间上的单调映射}\label{subsec:monotone-banach-mappings}

本节将看到, 可以把前一节的一些思想推广到无限维空间. 假设 $X$ 为实自反 Banach 空间, 研究映射 $T:X\to X^*$. 对元素 $g\in X^*$, 使用标准内积记号, 将 $g(v)$ 写成 $(g,v)$. \emph{在本章的其余各节中, 为使叙述更易理解, 只在 $X$ 和 $X^*$ 都可分的情形下给出证明.}

首先给出一系列定义, 将有限维空间上函数的概念推广到自反 Banach 空间上的映射.

\begin{Definition}
\label{def:banach-nonlinear-bounded-continuous-mapping}
若映射 $T:X\to X^*$ 将 $X$ 中的有界集映到 $X^*$ 中的有界集, 则称它为\textbf{有界的}. 若对每个 $u\in X$ 都有
\begin{equation}\label{eq:banach-nonlinear-mapping-continuity-definition}
\|T(u)-T(v)\|_{X^*}\to0\quad\text{whenever }\|u-v\|_X\to0.
\end{equation}
则称该映射为\textbf{连续的}.
\end{Definition}

\begin{Definition}
\label{def:banach-nonlinear-monotone-mapping}
若
\begin{equation}\label{eq:banach-nonlinear-mapping-monotonicity-definition}
(T(u)-T(v),u-v)\geq0\quad\text{for all }u,v\in X,
\end{equation}
则称映射 $T:X\to X^*$ 为单调的; 若每当 $u\neq v$ 时该不等式都严格成立, 则称它为严格单调的.
\end{Definition}

\begin{Definition}
\label{def:banach-nonlinear-coercive-mapping}
若
\begin{equation}\label{eq:banach-nonlinear-mapping-coercivity-definition}
\dfrac{(T(u),u)}{\|u\|}\to\infty\quad\text{as }\|u\|\to\infty.
\end{equation}
则称映射 $T:X\to X^*$ 为强制的.
\end{Definition}

\begin{Remark}
\label{rem:finite-banach-pairing-generalization}
利用上面的定义, 前一节的结果容易推广到从 $n$ 维 Banach 空间 $X$ 到另一个 $n$ 维 Banach 空间 $Y$ 的映射 $T$. 这里唯一真正的技巧是定义恰当的双线性型
\begin{equation}\label{eq:finite-banach-spaces-appropriate-bilinear-pairing}
X\times Y\ni x,y\mapsto\langle x,y\rangle\in\mathbb{R}
\end{equation}
以代替本节定义中的 ``内积\textquotesingle\textquotesingle
\begin{equation}\label{eq:banach-space-dual-inner-product-notation}
X\times X^*\ni x,y\mapsto(x,y)\in\mathbb{R}
\end{equation}
以及前一节中的点积. 细节留给读者.
\end{Remark}

Lemma~\ref{lem:finite-dimensional-minty-equivalence} 和 Theorem~\ref{thm:finite-dimensional-monotone-solution-set-closed-convex} 的证明基于单调性, 丝毫不依赖于所涉及空间的维数. 因此, 只需对前面的证明作少量记号上的修改, 就能得到下面两个类似结果.

\begin{Lemma}
\label{lem:banach-minty-equivalence}
给定 $g\in X^*$, 并设 $T:X\to X^*$ 连续. 若 $u\in X$ 为变分不等式
\begin{equation}\label{eq:banach-minty-variational-inequality}
(T(v)-g,v-u)\geq0\quad\text{for every }v\in X,
\end{equation}
的解, 则 $u$ 满足方程
\begin{equation}\label{eq:banach-minty-equation}
T(u)=g.
\end{equation}
此外, 若 $T$ 单调, 则方程 \eqref{eq:banach-minty-equation} 的每个解也都是变分不等式 \eqref{eq:banach-minty-variational-inequality} 的解.
\end{Lemma}

\begin{Theorem}
\label{thm:banach-monotone-solution-set-closed-convex}
设 $g\in X^*$, 并设 $T:X\to X^*$ 连续且单调. 则 $T(u)=g$ 的解集 $\mathcal{K}$ 为闭凸集.
\end{Theorem}

下面给出 Theorem~\ref{thm:finite-dimensional-coercive-surjectivity} 的第一个无限维类似结果.

\begin{Theorem}{Browder-Minty}
\label{thm:browder-minty-surjectivity}
设 $X$ 为实自反 Banach 空间, 并设 $T:X\to X^*$ 有界, 连续, 强制且单调. 则对任意 $g\in X^*$, 方程
\begin{equation}\label{eq:browder-minty-surjectivity-equation}
T(u)=g;
\end{equation}
存在解 $u$; 即 $T(X)=X^*$.
\end{Theorem}

\begingroup
\renewcommand{\endProof}{\par\nobreak\hfill$\blacksquare$\par\vspace{0.1cm}}
\begin{Proof}
\PDEIndexMark{idx:browder-minty-entire-galerkin-limit-proof}设 $\{x_i\}_{i=1}^\infty$ 为线性张成在 $X$ 中稠密的一个集合. (回顾, 这里只在 $X$ 和 $X^*$ 都可分的情形下给出证明.) 现在令
\begin{equation}\label{eq:browder-minty-dense-finite-dimensional-spaces}
X\supset X_n:=\operatorname{span}\{x_1,x_2,\ldots,x_n\}.
\end{equation}
计划用一个有限维问题
\begin{equation}\label{eq:browder-minty-galerkin-finite-dimensional-problem}
T_n(u_n)=g_n,
\end{equation}
逼近问题 \eqref{eq:browder-minty-surjectivity-equation}, 其中 $u_n\in X_n$, $g_n\in X_n^*$, 且 $T_n:X_n\to X_n^*$. 用有限维问题逼近无限维问题的思想称为 \emph{Galerkin 方法}. 它作为数值计算的手段广为人知, 但作为理论工具也同样有用 (正如这里的用法).

假设已给定 $g\in X^*$, 将 $g_n\in X_n^*$ 定义为限制
\begin{equation}\label{eq:browder-minty-restricted-right-hand-side}
g_n:=g|_{X_n}.
\end{equation}
类似地, 定义算子 $T_n:X_n\to X_n^*$ 为
\begin{equation}\label{eq:browder-minty-restricted-nonlinear-mapping}
T_n(u):=T(u)|_{X_n}\quad\text{for all }u\in X_n.
\end{equation}
另一种说法是, 对任意 $u,v\in X_n$, 有
\begin{equation}\label{eq:browder-minty-restriction-pairing-identities}
(g-g_n,v)=(T(u)-T_n(u),v)=0.
\end{equation}
$T_n$ 的连续性, 有界性和强制性直接继承自 $T$ 的性质. 因此, 利用 Theorem~\ref{thm:finite-dimensional-coercive-surjectivity}, 可知存在序列 $\{u_n\}_{n=1}^\infty$, 使
\begin{equation}\label{eq:browder-minty-galerkin-solution-sequence}
T_n(u_n)=g_n.
\end{equation}
现在利用 \eqref{eq:browder-minty-restriction-pairing-identities}, 注意到
\begin{equation}\label{eq:browder-minty-coercive-solution-norm-bound}
\begin{aligned}
\dfrac{(T(u_n),u_n)}{\|u_n\|}&=\dfrac{(T_n(u_n),u_n)}{\|u_n\|}=\dfrac{(g_n,u_n)}{\|u_n\|}\\
&=\dfrac{(g,u_n)}{\|u_n\|}\leq\dfrac{\|g\|\,\|u_n\|}{\|u_n\|}=\|g\|.
\end{aligned}
\end{equation}
于是, 由于 $T$ 为强制的, $\|u_n\|$ 必有界. 又由于 $T$ 有界, $\|T(u_n)\|$ 必有界. 因此, 利用弱紧性定理 Theorem~\ref{thm:alaoglu-weak-star-subsequence}, $X$ 的自反性以及 $X^*$ 的可分性, 可知存在 $u\in X$, $\widetilde g\in X^*$ 和子序列 (仍记为 $u_n$), 使
\begin{equation}\label{eq:browder-minty-galerkin-weak-solution-limit}
u_n\rightharpoonup u\quad\text{in }X
\end{equation}
且
\begin{equation}\label{eq:browder-minty-galerkin-weak-image-limit}
T(u_n)\rightharpoonup\widetilde g\quad\text{in }X^*.
\end{equation}
现在证明 $\widetilde g=g$ 且 $T(u)=g$.

再次利用 \eqref{eq:browder-minty-restriction-pairing-identities}, 注意到对每个基向量 $x_i$,
\begin{equation}\label{eq:browder-minty-identify-weak-image-limit}
(\widetilde g-g,x_i)=\lim_{n\to\infty}(T(u_n)-g,x_i)=0.
\end{equation}
因此, $\widetilde g=g$.

这里的关键结果是, 再次利用 \eqref{eq:browder-minty-restriction-pairing-identities}, 可以得到
\begin{equation}\label{eq:browder-minty-galerkin-weak-pairing-limit}
\begin{aligned}(T(u_n),u_n)&=(T_n(u_n),u_n)\\&=(g_n,u_n)\\&=(g,u_n)\\&\to(g,u).\end{aligned}
\end{equation}
(由于左端为两个\emph{弱}收敛序列的乘积, 这一结果并非平凡的.) 现在利用这一事实和 $T$ 的单调性, 得到对每个 $v\in X$ 都有
\begin{equation}\label{eq:browder-minty-monotonicity-limit-inequality}
0\leq(T(u_n)-T(v),u_n-v)\to(T(v)-g,v-u).
\end{equation}
由 Lemma~\ref{lem:banach-minty-equivalence} 立即得到定理的结论.
\end{Proof}
\endgroup

\subsubsection{单调算子对非线性 PDE 的应用}\label{subsec:monotone-operators-nonlinear-pde-applications}

本节将单调算子理论应用于散度形式的标量二阶拟线性 PDE. 因此, 假设 $\Omega\subset\mathbb{R}^n$ 为区域, 对充分光滑的函数 $u:\Omega\to\mathbb{R}$, 令
\begin{equation}\label{eq:monotone-quasilinear-divergence-operator}
A(u)=\sum_{|\alpha|\leq1}(-1)^{|\alpha|}D^\alpha\mathcal{A}_\alpha(\bm{x},\delta_1(u(\bm{x}))),
\end{equation}
其中
\begin{equation}\label{eq:monotone-quasilinear-first-derivative-jet}
\delta_1(u(\bm{x})):=\{D^\alpha u(\bm{x})\mid|\alpha|\leq1\}.
\end{equation}
将考虑方程
\begin{equation}\label{eq:monotone-quasilinear-dirichlet-equation}
A(u)=f
\end{equation}
的 Dirichlet 问题, 其中 $f\in W^{-1,q}(\Omega)$, 给定 $p\in(1,\infty)$, 且 $\frac{1}{p}+\frac{1}{q}=1$. (这里类似于 Section~\ref{sec:negative-sobolev-duality}, 将 $W^{-1,q}(\Omega)$ 定义为 $W_0^{1,p}(\Omega)$ 的对偶空间.) 对应于 $A$ 的 ``二元型\textquotesingle\textquotesingle\ 为
\begin{equation}\label{eq:monotone-quasilinear-bivariate-form}
B(u,v):=\sum_{|\alpha|\leq1}\int_\Omega\mathcal{A}_\alpha(\bm{x},\delta_1(u(\bm{x})))D^\alpha v(\bm{x})\,\mathrm{d}\bm{x}.
\end{equation}


\begin{Definition}
\label{def:monotone-quasilinear-dirichlet-weak-solution}
若
\begin{equation}\label{eq:monotone-quasilinear-dirichlet-weak-formulation}
B(u,v)=(f,v)\quad\text{for all }v\in W_0^{1,p}(\Omega).
\end{equation}
则称 $u\in W_0^{1,p}(\Omega)$ 为拟线性 PDE \eqref{eq:monotone-quasilinear-dirichlet-equation} 的 Dirichlet 问题的\textbf{弱解}.
\end{Definition}

为得到存在性定理, 对函数 $\mathcal{A}_\alpha$ 作如下假设.

\begin{enumerate}[label=H-\arabic*.,ref=H-\arabic*]\item\label{hyp:monotone-pde-measurability} 对每个 $|\alpha|\leq1$,
\begin{equation}\label{eq:monotone-pde-coefficient-measurability}
\bm{x}\mapsto\mathcal{A}_\alpha(\bm{x},\delta_1)
\end{equation}
对每个固定的 $\delta_1=(\bm{x}i_\alpha)_{|\alpha|\leq1}\in\mathbb{R}^{n+1}$ 都是可测的.

\item\label{hyp:monotone-pde-continuity} 对每个 $|\alpha|\leq1$,
\begin{equation}\label{eq:monotone-pde-coefficient-jet-continuity}
\delta_1\mapsto\mathcal{A}_\alpha(\bm{x},\delta_1)
\end{equation}
对几乎每个 $\bm{x}\in\Omega$ 都属于 $C(\mathbb{R}^{n+1})$.

\item\label{hyp:monotone-pde-monotonicity} 对每个 $\delta_1^1=(\bm{x}i_\alpha^1)_{|\alpha|\leq1}\in\mathbb{R}^{n+1}$, $\delta_1^2=(\bm{x}i_\alpha^2)_{|\alpha|\leq1}\in\mathbb{R}^{n+1}$ 以及每个 $\bm{x}\in\Omega$, 都有
\begin{equation}\label{eq:monotone-pde-coefficient-monotonicity}
\sum_{|\alpha|\leq1}[\mathcal{A}_\alpha(\bm{x},\delta_1^1)-\mathcal{A}_\alpha(\bm{x},\delta_1^2)](\bm{x}i_\alpha^1-\bm{x}i_\alpha^2)\geq0.
\end{equation}


\item\label{hyp:monotone-pde-coercivity} 存在 $p\in(1,\infty)$, 常数 $c_0>0$ 和函数 $h\in L^1(\Omega)$, 使对每个 $\bm{x}\in\Omega$ 以及每个 $\delta_1=(\bm{x}i_\alpha)_{|\alpha|\leq1}\in\mathbb{R}^{n+1}$, 都有
\begin{equation}\label{eq:monotone-pde-coefficient-coercivity}
\sum_{|\alpha|\leq1}\mathcal{A}_\alpha(\bm{x},\delta_1)\bm{x}i_\alpha\geq c_0|\delta_1|^p-h(\bm{x}).
\end{equation}


\item\label{hyp:monotone-pde-growth} 存在常数 $c_1>0$ 和函数 $g\in L^q(\Omega)$, ($q:=p/(p-1)$), 使对每个 $\bm{x}\in\Omega$ 以及每个 $\delta_1=(\bm{x}i_\alpha)_{|\alpha|\leq1}\in\mathbb{R}^{n+1}$, 都有
\begin{equation}\label{eq:monotone-pde-coefficient-growth-bound}
|\mathcal{A}_\alpha(\bm{x},\delta_1)|\leq c_1|\delta_1|^{p-1}+g(\bm{x}).
\end{equation}
\end{enumerate}

在这些假设下, 得到下面的结果.

\begin{Theorem}
\label{thm:monotone-quasilinear-dirichlet-existence}
设 $\mathcal{A}_\alpha$ 满足假设 \ref{hyp:monotone-pde-measurability}--\ref{hyp:monotone-pde-growth}, 并设 $p$ 和 $q$ 分别如假设 \ref{hyp:monotone-pde-coercivity} 和 \ref{hyp:monotone-pde-growth} 中所定义. 则对每个 $f\in W^{-1,q}(\Omega)$, 拟线性 PDE \eqref{eq:monotone-quasilinear-dirichlet-equation} 的 Dirichlet 问题存在弱解 $u\in W_0^{1,p}(\Omega)$.
\end{Theorem}

将这一定理的证明分成一系列引理, 第一个如下.

\begin{Lemma}
\label{lem:monotone-pde-induced-bounded-dual-mapping}
对每个固定的 $u\in W_0^{1,p}(\Omega)$, 映射
\begin{equation}\label{eq:monotone-pde-bounded-test-functional}
W_0^{1,p}(\Omega)\ni v\mapsto B(u,v)\in\mathbb{R}
\end{equation}
为有界线性泛函. 因此, 存在映射
\begin{equation}\label{eq:monotone-pde-induced-dual-mapping}
X=W_0^{1,p}(\Omega)\ni u\mapsto T(u)\in W^{-1,q}(\Omega)=X^*
\end{equation}
使得
\begin{equation}\label{eq:monotone-pde-induced-duality-pairing}
B(u,v)=(T(u),v),
\end{equation}
对所有属于 $W_0^{1,p}(\Omega)$ 的 $u,v$ 成立. 此外, 该非线性映射有界.
\end{Lemma}

\begingroup
\renewcommand{\endProof}{\par\nobreak\hfill$\blacksquare$\par\vspace{0.1cm}}
\begin{Proof}
映射 $v\mapsto B(u,v)$ 的线性是显然的. 若 $u,v\in W_0^{1,p}(\Omega)$, 则利用 \ref{hyp:monotone-pde-growth} 和 Hölder 不等式, 得到估计:\[\begin{aligned}
|B(u,v)|&\leq\sum_{|\alpha|\leq1}\int_\Omega|\mathcal{A}_\alpha(\bm{x},\delta_1(u(\bm{x})))D^\alpha v(\bm{x})|\,\mathrm{d}\bm{x}\\
&\leq\sum_{|\alpha|\leq1}\int_\Omega[c_1|\delta_1(u(\bm{x}))|^{p-1}+g(\bm{x})]|D^\alpha v(\bm{x})|\,\mathrm{d}\bm{x}\\
&\leq\sum_{|\alpha|\leq1}c_1\left\{\int_\Omega|\delta_1(u(\bm{x}))|^{q(p-1)}\right\}^{1/q}\left\{\int_\Omega|D^\alpha v|^p\,\mathrm{d}\bm{x}\right\}^{1/p}\\
&\quad+\sum_{|\alpha|\leq1}\left\{\int_\Omega|g|^q\,\mathrm{d}\bm{x}\right\}^{1/q}\left\{\int_\Omega|D^\alpha v|^p\,\mathrm{d}\bm{x}\right\}^{1/p}\\
&\leq C(\|u\|_{1,p}^{p-1}+1)\|v\|_{1,p}.
\end{aligned}\]这表明对每个 $u\in W_0^{1,p}(\Omega)$, 映射 $v\mapsto B(u,v)$ 都有界. 这就蕴含映射 $T$ 的存在性.

此外, 同一估计还表明 $T$ 有界, 因为
\begin{equation}\label{eq:monotone-pde-induced-dual-norm-bound}
\|T(u)\|_{-1,q}=\sup_{\substack{v\in W_0^{1,p}(\Omega)\\\|v\|_{1,p}=1}}|(T(u),v)|\leq C(\|u\|_{1,p}^{p-1}+1).
\end{equation}
因此, $T$ 将 $W_0^{1,p}(\Omega)$ 中的有界集映到 $W^{-1,q}(\Omega)$ 中的有界集. 证明完成.
\end{Proof}
\endgroup

\begin{Lemma}
\label{lem:monotone-pde-induced-mapping-monotone-coercive}
Lemma~\ref{lem:monotone-pde-induced-bounded-dual-mapping} 中定义的算子 $T$ 为单调且强制的.
\end{Lemma}

\begingroup
\renewcommand{\endProof}{\par\nobreak\hfill$\blacksquare$\par\vspace{0.1cm}}
\begin{Proof}
为证明单调性, 利用假设 \ref{hyp:monotone-pde-monotonicity}, 得到\[\begin{aligned}
(T(u)-T(v),u-v)&=B(u,u-v)-B(v,u-v)\\
&=\sum_{|\alpha|\leq1}\int_\Omega[\mathcal{A}_\alpha(\bm{x},\delta_1(u))-\mathcal{A}_\alpha(\bm{x},\delta_1(v))][D^\alpha u-D^\alpha v]\,\mathrm{d}\bm{x}\\
&\geq0.
\end{aligned}\]为证明强制性, 利用假设 \ref{hyp:monotone-pde-coercivity}, 得到\[\begin{aligned}
\dfrac{(T(u),u)}{\|u\|_{1,p}}&=\dfrac{B(u,u)}{\|u\|_{1,p}}\\
&=\|u\|_{1,p}^{-1}\sum_{|\alpha|\leq1}\int_\Omega\mathcal{A}_\alpha(\bm{x},\delta_1(u))D^\alpha u\,\mathrm{d}\bm{x}\\
&\geq\|u\|_{1,p}^{-1}\left[c_0\int_\Omega|\delta_1(u)|^p\,\mathrm{d}\bm{x}-\int_\Omega h(\bm{x})\,\mathrm{d}\bm{x}\right]\\
&\geq\|u\|_{1,p}^{-1}C[\|u\|_{1,p}^p-1].
\end{aligned}\]由于 $p>1$, 有
\begin{equation}\label{eq:monotone-pde-induced-mapping-coercivity-limit}
\dfrac{(T(u),u)}{\|u\|_{1,p}}\to\infty\quad\text{as }\|u\|_{1,p}\to\infty.
\end{equation}
证明完成.
\end{Proof}
\endgroup

因此, 为将 Browder-Minty 定理应用于映射 $T$, 并完成 Theorem~\ref{thm:monotone-quasilinear-dirichlet-existence} 的证明, 只需证明下面的结论.

\begin{Lemma}
\label{lem:monotone-pde-induced-mapping-continuity}
映射 $T:W_0^{1,p}(\Omega)\to W^{-1,q}(\Omega)$ 连续.
\end{Lemma}

下一节将介绍一种称为 \emph{Nemytskii 算子}的工具, 用它可证明这一引理.

\begin{Example}
\label{ex:anisotropic-p-power-nonlinear-elliptic-operator}
考虑二阶非线性偏微分算子
\begin{equation}\label{eq:anisotropic-p-power-divergence-example}
A(u):=-\sum_{i=1}^n\dfrac{\partial}{\partial x_i}\left(\left|\dfrac{\partial u}{\partial x_i}\right|^{p-2}\dfrac{\partial u}{\partial x_i}\right)+|u|^{p-2}u,
\end{equation}
其中 $p\in(1,\infty)$. 注意, $p=2$ 时就是 Laplacian 加一个低阶项, 这已在线性问题的内容中讨论过. 这里,\[\begin{aligned}
\mathcal{A}_{(0,0,\ldots,0)}&=a_0(\bm{x},u,u_{x_1},\ldots,u_{x_n})=|u|^{p-2}u,\\
\mathcal{A}_{(0,\ldots,0,1,0,\ldots,0)}&=a_i(\bm{x},u,u_{x_1},\ldots,u_{x_n})=|u_{x_i}|^{p-2}u_{x_i},\quad i=1,\ldots,n.
\end{aligned}\]要验证这些 $\mathcal{A}_\alpha$ 满足假设 \ref{hyp:monotone-pde-measurability} 到 \ref{hyp:monotone-pde-growth}. 假设 \ref{hyp:monotone-pde-measurability} 和 \ref{hyp:monotone-pde-continuity} 显然成立. 为验证 \ref{hyp:monotone-pde-monotonicity}, 令 $\delta_1^1=(\bm{x}i_0^1,\bm{x}i_1^1,\ldots,\bm{x}i_n^1)$ 和 $\delta_1^2=(\bm{x}i_0^2,\bm{x}i_1^2,\ldots,\bm{x}i_n^2)$, 计算得到\[\begin{aligned}
&\sum_{|\alpha|\leq1}[\mathcal{A}_\alpha(\bm{x},\delta_1^1)-\mathcal{A}_\alpha(\bm{x},\delta_1^2)](\bm{x}i_\alpha^1-\bm{x}i_\alpha^2)\\
&\quad=\sum_{i=0}^n\bigl(|\bm{x}i_i^1|^{p-2}\bm{x}i_i^1-|\bm{x}i_i^2|^{p-2}\bm{x}i_i^2\bigr)(\bm{x}i_i^1-\bm{x}i_i^2)\\
&\quad\geq0.
\end{aligned}\]为验证 \ref{hyp:monotone-pde-coercivity}, 令 $\delta_1=(\bm{x}i_0,\bm{x}i_1,\ldots,\bm{x}i_n)$, 得到
\begin{equation}\label{eq:anisotropic-p-power-printed-coercivity-chain}
\sum_{|\alpha|\leq1}\mathcal{A}_\alpha(\bm{x},\delta_1)\bm{x}i_\alpha=\sum_{i=0}^n|\bm{x}i_i|^{p-2}\bm{x}i_i\bm{x}i_i=\sum_{i=1}^n|\bm{x}i_i|^p\geq c_0|\delta_i|^p
\end{equation}
由下式可知 \ref{hyp:monotone-pde-growth} 成立
\begin{equation}\label{eq:anisotropic-p-power-coefficient-growth}
|\mathcal{A}_\alpha(\bm{x},\delta_1)|=|a_i(\bm{x},\delta_1)|=|\bm{x}i_i|^{p-1}\leq|\delta_1|^{p-1}.
\end{equation}
因此, 下面的存在性结果立即由 Theorem~\ref{thm:monotone-quasilinear-dirichlet-existence} 得出.
\end{Example}

\begin{Theorem}
\label{thm:anisotropic-p-power-dirichlet-existence}
设非线性二阶偏微分算子 $A$ 由 \eqref{eq:anisotropic-p-power-divergence-example} 定义. 则对每个 $f\in W^{-1,q}(\Omega)$, 方程
\begin{equation}\label{eq:anisotropic-p-power-dirichlet-existence-equation}
A(u)=f.
\end{equation}
存在弱解 $u\in W_0^{1,p}$.
\end{Theorem}

\subsubsection{Nemytskii 算子}\label{subsec:nemytskii-operators}

本节将陈述一些关于 $L^p(\Omega)$ 与非线性函数复合的重要结果, 不给出证明. 更详细的论述可参阅 \cite{Li}.

\begin{Definition}
\label{def:caratheodory-conditions}
设 $\Omega\subset\mathbb{R}^n$ 为区域. 若函数
\begin{equation}\label{eq:caratheodory-nonlinear-function-domain}
\Omega\times\mathbb{R}^m\ni(\bm{x},\bm{u})\mapsto f(\bm{x},\bm{u})\in\mathbb{R}
\end{equation}
满足
\begin{equation}\label{eq:caratheodory-value-variable-continuity}
\bm{u}\mapsto f(\bm{x},\bm{u})\text{ is continuous for almost every }\bm{x}\in\Omega
\end{equation}
和
\begin{equation}\label{eq:caratheodory-spatial-measurability-printed-domain}
\bm{x}\mapsto f(\bm{x},\bm{u})\text{ is measurable for every }\bm{u}\in\Omega.
\end{equation}
则称它满足 \textbf{Carathéodory 条件}.
\end{Definition}

任给满足 Carathéodory 条件的 $f$ 以及函数 $\bm{u}:\Omega\to\mathbb{R}^m$, 可以通过复合定义另一个函数
\begin{equation}\label{eq:nemytskii-pointwise-composition-operator}
\mathcal{F}(\bm{u})(\bm{x}):=f(\bm{x},\bm{u}(\bm{x})).
\end{equation}
复合算子 $\mathcal{F}$ 称为 \emph{Nemytskii 算子}. 主要定理讨论这些算子从 $L^p(\Omega)$ 到 $L^q(\Omega)$ 的有界性和连续性.

\begin{Theorem}
\label{thm:nemytskii-lp-lq-bounded-continuous}
设 $\Omega\subset\mathbb{R}^n$ 为区域, 并设
\begin{equation}\label{eq:nemytskii-growth-theorem-function-domain}
\Omega\times\mathbb{R}^m\ni(\bm{x},\bm{u})\mapsto f(\bm{x},\bm{u})\in\mathbb{R}
\end{equation}
满足 Carathéodory 条件. 此外, 给定 $p\in(1,\infty)$ 和 $g\in L^q(\Omega)$ (其中 $\frac{1}{p}+\frac{1}{q}=1$), 并设 $f$ 满足
\begin{equation}\label{eq:nemytskii-growth-bound-lp-lq}
|f(\bm{x},\bm{u})|\leq C|\bm{u}|^{p-1}+g(\bm{x}).
\end{equation}
则由 \eqref{eq:nemytskii-pointwise-composition-operator} 定义的 Nemytskii 算子 $\mathcal{F}$ 为从 $L^p(\Omega)$ 到 $L^q(\Omega)$ 的有界连续映射.
\end{Theorem}

\begin{Remark}
\label{rem:nemytskii-monotone-pde-continuity-corollary}
Lemma~\ref{lem:monotone-pde-induced-mapping-continuity} 是这一定理的推论. 为看出这一点, 只需应用假设 \ref{hyp:monotone-pde-measurability}, \ref{hyp:monotone-pde-continuity} 和 \ref{hyp:monotone-pde-growth}, 就可知每个 $\mathcal{A}_\alpha$ 都可以作为满足适当增长条件的 Nemytskii 算子. $T$ 从 $W^{1,p}(\Omega)$ 到 $W^{-1,q}(\Omega)$ 的连续性, 由 $\delta_1(\bm{x})\mapsto\mathcal{A}_\alpha(\bm{x},\delta_1(\bm{x}))$ 作为从 $L^p(\Omega)$ 到 $L^q(\Omega)$ 的映射的连续性得到.
\end{Remark}
